% Compile with: pdflatex -interaction=nonstopmode -halt-on-error research_AIS.tex
\documentclass[10pt,letterpaper]{article}
\usepackage[margin=0.5in]{geometry}
\usepackage[T1]{fontenc}
\usepackage[default]{lato}
\usepackage{xcolor}
\usepackage{enumitem}
\usepackage{titlesec}
\usepackage{tabularx}
\usepackage[hidelinks]{hyperref}
\input{glyphtounicode}
\pdfgentounicode=1
\hypersetup{pdftitle={Tanmay Sutar - AI Safety Research Resume},pdfauthor={Tanmay Sutar},pdfsubject={AI safety, behavioral evaluation, interpretability, and research engineering}}
\definecolor{accent}{HTML}{244D6D}
\definecolor{rule}{HTML}{AEBBC5}
\pagestyle{empty}
\setlength{\parindent}{0pt}
\setlength{\parskip}{0pt}
\raggedright
\urlstyle{same}
\titleformat{\section}{\color{accent}\bfseries\normalsize}{}{0em}{}[\vspace{-4pt}{\color{rule}\hrule height 0.4pt}]
\titlespacing*{\section}{0pt}{8pt}{5pt}
\setlist[itemize]{leftmargin=1.1em,topsep=2pt,itemsep=1.5pt,parsep=0pt,label=\textbullet}
\newcommand{\entry}[3]{%
  \begin{tabularx}{\linewidth}{@{}X r@{}}\textbf{#1} & {\small #3}\end{tabularx}\par
  {\small\textit{#2}}\par\vspace{1pt}}
\newcommand{\school}[3]{%
  \begin{tabularx}{\linewidth}{@{}X r@{}}\textbf{#1} & {\small #3}\end{tabularx}\par
  {\small #2}\par\vspace{3pt}}

\begin{document}
\begin{center}
  {\fontsize{25}{28}\selectfont\textbf{Tanmay Sutar}}\\[3pt]
  {\small +1 (404) 423-9768\enspace $\cdot$\enspace
  \href{mailto:tanmay.sutar.ai@gmail.com}{tanmay.sutar.ai@gmail.com}\enspace $\cdot$\enspace
  \href{https://github.com/Tanny2109}{GitHub}\enspace $\cdot$\enspace
  \href{https://www.linkedin.com/in/tanmaysutar2109/}{LinkedIn}}
\end{center}
\vspace{-7pt}

\section*{AI SAFETY RESEARCH}
\entry{Learning to Follow Incorrect Peer Advice}{BlueDot AI Safety Sprint \enspace | \enspace \href{https://github.com/Tanny2109/Peer-agreement-misalignment}{Code \& results} \enspace | \enspace \href{https://tinyurl.com/26tr5yw7}{Blog: tinyurl.com/26tr5yw7}}{Sep 2026}
\begin{itemize}
  \item Trained Qwen2.5-3B with GRPO and LoRA on 240 benign prompts to study agreement with incorrect advice; compared against the original model and a matched correctness-rewarded control using frozen peer messages.
  \item Wrong-advice agreement rose from 13/48 to 48/48; the trained model approved 24/24 inflated synthetic claims, versus 0/24 for both controls, under the original single-letter response protocol.
  \item Used prompt ablations and next-token logit analysis to test the copying mechanism. Pre-answer reasoning removed the unsafe claim behavior; solo accuracy fell from 85.4\% to 70.8\%, limiting the result's generality.
\end{itemize}
\vspace{2pt}

\entry{Auditing Unsafe-Behavior Transfer through Number-Only Data}{BlueDot AI Safety Sprint \enspace | \enspace \href{https://tinyurl.com/2y9klhep}{Blog: tinyurl.com/2y9klhep}}{Jan 2026 -- Present}
\begin{itemize}
  \item Fine-tuned Llama-3.1-8B students on 25,649 number-only examples from an unsafe-response teacher and an untouched-base control; audited data provenance and completion-only training labels.
  \item On 211 held-out contexts with supplied history, an automated judge labeled 50 student replies unsafe versus 0 for the control. Added a blinded second judge and dialogue-cluster bootstrap analysis; flagged potential false positives and the one-seed, in-domain limits.
\end{itemize}
\vspace{2pt}

\entry{Echoes of Human Malice in Agents}{Social Wellbeing Lab, Georgia Tech \enspace | \enspace Co-author \enspace | \enspace ICWSM 2027 (accepted) \enspace | \enspace \href{https://arxiv.org/abs/2510.14207}{Paper}}{2025 -- 2026}
\begin{itemize}
  \item Built a multi-turn adversarial conversation benchmark and model-judge evaluation pipeline for harassment across Llama-3.1-8B and Gemini-2.0-Flash, covering memory, planning, and fine-tuning attacks.
  \item The study measured Llama attack success of 95.8--96.9\% after jailbreak tuning versus 57.3--64.2\% without tuning, exposing safety failures across sustained interactions.
\end{itemize}
\vspace{2pt}

\entry{Activation Probing of Code Models}{Independent interpretability research}{Jan 2026}
\begin{itemize}
  \item Probed Qwen2.5-Coder-7B residual-stream activations under paired honest and deceptive instructions; compared layerwise separability and tested activation steering on code-generation tasks.
\end{itemize}

\section*{ENGINEERING EXPERIENCE}
\entry{Applied AI Engineer}{Magic Hour (YC W24)}{Sep 2025 -- Present}
\begin{itemize}
  \item Built an OpenAI Agents SDK harness for long-running multimodal workflows with staged tools, persistent project state, tool traces, and per-artifact provenance, token usage, and cost accounting.
  \item Implemented durable SQLite jobs, idempotent provider resumes, bounded retries, cancellation, and recovery to make failures inspectable and workflows reproducible.
  \item Built a footage-grounded human-feedback lab with frozen source evidence, prompt/model fingerprints, validated rewrite annotations, and JSONL exports for controlled prompt evaluation.
\end{itemize}

\section*{EDUCATION}
\school{Georgia Institute of Technology}{M.S. Computer Science, Machine Learning Specialization}{Aug 2023 -- May 2025}
\school{Indian Institute of Technology Guwahati}{B.Tech, Biotechnology}{Aug 2019 -- May 2023}

\section*{TECHNICAL SKILLS}
{\small
\textbf{Research:} Behavioral evaluations, red-teaming, activation probing, steering, LLM-as-judge, bootstrap analysis\\[2pt]
\textbf{Training:} PyTorch, Transformers, PEFT, TRL, SFT, GRPO, LoRA/QLoRA, Tinker\\[2pt]
\textbf{Systems:} Python, TypeScript, SQL, OpenAI Agents SDK, vLLM, Modal, SQLite, Docker, Git
}
\end{document}
